% !TeX TS-program = xelatex
% 编译流程：xelatex → makeindex `gind.ist` → 两遍 xelatex
% makeindex -s gind.ist mdframed-doc.idx



\setcounter{errorcontextlines}{999}
\documentclass[parskip=false,11pt,]{ltxmdf}
\usepackage{pdfpages}
\usepackage[framemethod=TikZ]{mdframed}
\usetikzlibrary{shadows}
\surroundwithmdframed[middlelinecolor=ltxmdfblue,middlelinewidth=1pt,%
                      roundcorner=10pt,innertopmargin=0pt,%
                      leftmargin=1cm,rightmargin=1cm,%
                      innerleftmargin=-15pt,innerrightmargin=-15pt,%
                      ignorelastdescenders,%
                      settings={\lstset{resetmargins}},%
                      skipbelow=\topskip,skipabove=\topskip,%
                      innerbottommargin=0pt,backgroundcolor=gray!10]%
                     {tltxmdfexample}

\EnableCrossrefs
\CodelineIndex
\OnlyDescription
\begin{document}

 \CharacterTable
  {Upper-case    \A\B\C\D\E\F\G\H\I\J\K\L\M\N\O\P\Q\R\S\T\U\V\W\X\Y\Z
   Lower-case    \a\b\c\d\e\f\g\h\i\j\k\l\m\n\o\p\q\r\s\t\u\v\w\x\y\z
   Digits        \0\1\2\3\4\5\6\7\8\9
   Exclamation   \!     Double quote  \"     Hash (number) \#
   Dollar        \$     Percent       \%     Ampersand     \&
   Acute accent  \'     Left paren    \(     Right paren   \)
   Asterisk      \*     Plus          \+     Comma         \,
   Minus         \-     Point         \.     Solidus       \/
   Colon         \:     Semicolon     \;     Less than     \<
   Equals        \=     Greater than  \>     Question mark \?
   Commercial at \@     Left bracket  \[     Backslash     \\
   Right bracket \]     Circumflex    \^     Underscore    \_
   Grave accent  \`     Left brace    \{     Vertical bar  \|
   Right brace   \}     Tilde         \~}

 \GetFileInfo{mdframed.sty}


 \title{\Pack{mdframed} 宏包
       \footnote{对宏包 \texttt{framed.sty} 的扩展}}
 \subtitle{自动跨页的边框环境}
 \author{\href{mailto:marco.daniel@mada-nada.de}{Marco Daniel},
        \href{mailto:marco.daniel@mada-nada.de}{Elke Schubert}}
 \version{\mdversion}
 \date{\mdfmaindate}
 \introduction{给文字加框的常规方法（\Cmd{fbox}
              或 \Cmd{fcolorbox}）需要手动处理分页，也就是说
              你不得不把 \Cmd{fbox} 拆成两个。本宏包定义了环境
              \Pack{mdframed}，能够自动处理带框文字的分页。\par
              通过定义新的环境，用户可以在多种个性化的
              外观之间进行选择。%
             \par\kern.5\baselineskip\noindent
             相关文件：\quad
              \begin{tabular}[t]{ll}
               \url{mdframed-example-default.pdf} &
               \url{mdframed-example-tikz.pdf} \\
               \url{mdframed-example-pstricks.pdf}  &
               \url{mdframed-example-texsx.pdf}
              \end{tabular}
             \par\kern.5\baselineskip\noindent
             备注：\quad\parbox[t]{.8\linewidth}{%
                我在 \href{https://github.com/marcodaniel/mdframed}{GitHub}
                上为 \Pack{mdframed} 建立了代码仓库，你可以在那里
                \href{https://github.com/marcodaniel/mdframed}{下载}
                当前的开发版本。}%
 }

 \maketitle

 \vspace*{2\baselineskip} 

 \section{设计动机}
 许多用户希望（进一步）突出引理、定义、证明等内容。
  \Pack{mdframed} 宏包允许你创建带可跨页边框的环境。
 我认为示例是展示其特性的最佳方式。
 
 \newmdtheoremenv[%
    outerlinewidth  = 2             ,%
    roundcorner     = 10pt          ,%
    leftmargin      = 40            ,%
    rightmargin     = 40            ,%
    backgroundcolor = yellow!40     ,%
    outerlinecolor  = blue!70!black ,% 
    innertopmargin  = \topskip      ,%
    splittopskip    = \topskip      ,%
    ntheorem        = true          ,% 
   ]{theorem}{定理}[section]
 
 \begin{theorem}[勾股定理]
  在任意直角三角形中，以斜边为边长的正方形面积，
  等于以两条直角边为边长的正方形面积之和。
  
 \[ a^2+b^2=c^2 \]
 \end{theorem}
 
 该边框由以下设置定义。
 
\begin{tltxmdfexample}[columns=fixed,fontadjust=false]
 \newmdtheoremenv[%
    outerlinewidth  = 2             ,%
    roundcorner     = 10pt          ,%
    leftmargin      = 40            ,%
    rightmargin     = 40            ,%
    backgroundcolor = yellow!40     ,%
    outerlinecolor  = blue!70!black ,% 
    innertopmargin  = \topskip      ,%
    splittopskip    = \topskip      ,%
    ntheorem        = true          ,% 
   ]{theorem}{定理}[section]
 \begin{theorem}[勾股定理]
 ...
 \end{theorem}
\end{tltxmdfexample}
 
 
 \section{语法}\label{sec:syntax}
 
 \minisec{\Pack{mdframed} 所需的宏包}
 宏包本身会加载以下宏包：
  \begin{center}
   \Pack{kvoptions}\qquad\Pack{xparse}\qquad\Pack{etoolbox}
   \quad\Pack{color}.
  \end{center}
 根据选项 \Opt{framemthod} 的不同，\Pack{mdframed} 还会加载：
  \begin{center}
   \Pack{xcolor}\qquad\Pack{tikz}\qquad\Pack{pstricks}.
  \end{center}
 
 按常规方式加载本宏包：
\begin{tltxmdfexample}
  \usepackage[<全局选项>]{mdframed}
\end{tltxmdfexample}
 
 只有选项 \Opt{framemethod} 应当通过 \Cmd{usepackage} 的
 可选参数传入；其余选项都应通过 \Cmd{mdfsetup} 或相关环境设置。
 如果你需要用到 \Pack{amsthm}，应在它之后再加载本宏包。
 \minisec{提供的环境}
 本宏包只定义了一个环境，语法如下：
\begin{tltxmdfexample}
  \begin{mdframed}[<局部选项>]
     <内容>
  \end{mdframed}
\end{tltxmdfexample}
 关于如何用 \Pack{mdframed} 创建自定义环境，请参见 \ref{sec:cmds} 节。
 
 \minisec{浮动体的自动检测}
 \Pack{mdframed} 会检测当前环境是否处于 \Pack{float}
 或 \Pack{minipage} 环境内部。若在这类环境中使用
 \Pack{mdframed}，它会自动启用 \Opt{nobreak} 选项。
 
 
 
 %%%%%%%%%%%%%%%%%%%%%%%
 \section{边框}\label{sec:frames}
 
 通常来说，\Pack{mdframed} 只会绘制若干线条。
 为支持跨页，宏包提供了以下几种边框样式。
 若以 \Opt{framemethod=default} 加载本宏包，
 则只能绘制单一线条。下图灰色方框内部即为将要排版文字的区域。
 \newcommand\Punkt[2][below]{%
 	\node[label=#1:#2] at (#2){\tikz\fill(0,0)circle(0.05);};}
 \newcommand\Single[1][]{\path[#1] (O)--(O|-P)--(P)--(P|-O)--cycle;}
 \newcommand\FirstSecond[1][]{\path[#1] (O|-P)--(O)--(P|-O)--(P);}
 \newcommand\MiddleFlaeche[1][]{\path[#1] (O)--(O|-P)--(P)--(P|-O);}
 \newcommand\MiddleLinien[1][]{\path[#1] (O)--(O|-P)(P)--(P|-O);}
 \begin{center}
 \vspace*{2\baselineskip}
 \null\hfill
 \begin{tikzpicture}[baseline=(current bounding box.north)]
 	\coordinate (O) at (0,0);
 	\pgfmathsetlengthmacro\x{4cm+2mm+6mm}
 	\pgfmathsetlengthmacro\y{2cm+2mm+6mm}
 	\coordinate (P) at (\x,\y);
 	\begin{scope}
 		\Single[line width=3mm,draw=blue!20]
 		\Single[clip]
 		\Single[fill=gray!70]
 		\Single[line width=7mm,draw=green!20]
 	\end{scope}
 	\Single[line width=2mm,draw=red!20,use as bounding box]
 	\node[anchor=south,font=\bfseries,text width=4cm,align=center] %
        at (24mm,35mm) {单页边框};
 \end{tikzpicture}%
 \hfill%
 \begin{tikzpicture}[baseline=(current bounding box.north)]
 	\coordinate (O) at (0,0);
 	\pgfmathsetlengthmacro\x{4cm+2mm+6mm}
 	\pgfmathsetlengthmacro\y{-2cm-1mm-3mm}
 	\coordinate (P) at (\x,\y);
 	\begin{scope}
 		\FirstSecond[line width=3mm,draw=blue!20]
 		\FirstSecond[clip]
 		\FirstSecond[fill=gray!70]
 		\FirstSecond[line width=7mm,draw=green!20]
 	\end{scope}
 	\FirstSecond[line width=2mm,draw=red!20,use as bounding box]
 	\node[anchor=south,font=\bfseries,text width=4cm,align=center] 
            at (2.4,0.7) {第一个框 —— 输出的第一部分};
 \end{tikzpicture}
 \hfill\null
 
 \vspace{3\baselineskip}
 
  \null\hfill
 \begin{tikzpicture}[baseline=(current bounding box.north)]
 	\coordinate (O) at (0,0);
 	\pgfmathsetlengthmacro\x{4cm+2mm+6mm}
 	\pgfmathsetlengthmacro\y{2cm+1mm+3mm}
 	\coordinate (P) at (\x,\y);
 	\begin{scope}
 		\FirstSecond[line width=3mm,draw=blue!20]
 		\FirstSecond[clip]
 		\FirstSecond[fill=gray!70]
 		\FirstSecond[line width=7mm,draw=green!20]
 	\end{scope}
 	\FirstSecond[line width=2mm,draw=red!20,use as bounding box]
 	\node[anchor=south,font=\bfseries,text width=4cm,align=center] %
        at (24mm,25mm) {第二个框 —— 输出的最后一部分};
 \end{tikzpicture}%
 \hfill%
 \begin{tikzpicture}[baseline=(current bounding box.north)]
 	\coordinate (O) at (0,0);
 	\pgfmathsetlengthmacro\x{4cm+2mm+6mm}
 	\pgfmathsetlengthmacro\y{2cm}
 	\coordinate (P) at (\x,\y);
 	\begin{scope}
 		\MiddleLinien[line width=3mm,draw=blue!20]
 		\MiddleFlaeche[clip]
 		\MiddleFlaeche[fill=gray!70]
 		\MiddleLinien[line width=7mm,draw=green!20]
 	\end{scope}
 	\MiddleLinien[line width=2mm,draw=red!20,use as bounding box]
 	\node[anchor=south,font=\bfseries,text width=4cm,align=center] 
        at (24mm,20mm) {中间的框 —— 输出的中间部分};
 \end{tikzpicture}
 \hfill\null
 
 \captionof{figure}{基本边框}\label{fig:basicframes}
\end{center} 
 %%%%%%%%%%%%%%%%%%%%%%%
 \section{命令}\label{sec:cmds}

 以下命令可在使用 \Pack{mdframed} 时为你提供帮助。
 
 \ExplCmd{newmdenv[选项]\{环境名\}} 该命令用于定义一个新环境，
 并由 \Env{mdframed} 为其加框。命令语法如下：
\begin{tltxmdfexample}
 \newmdenv[<mdframed 选项>]{环境名称}
\end{tltxmdfexample}
 这样就可以写出：
\begin{tltxmdfexample}
  \newmdenv[linecolor=red,frametitle=信息框]{infobox}

  \begin{infobox}[backgroundcolor=yellow]
   一些信息\ldots
  \end{infobox}
\end{tltxmdfexample}
 \ExplCmd{renewmdenv[选项]\{环境名\}} 
      使用该命令可以重新定义由 \Cmd{newmdenv}
      创建的环境。
 
 \ExplCmd{surroundwithmdframed[选项]\{环境\}} 
      有时你已经有预先定义好的环境。该命令可以在不改变
      原名称的前提下，用 \Env{mdframed} 给某个已有
      \Env{环境} 加框。例如要给 \Env{verbatim}
      环境加上 \Env{mdframed}，只需这样写：
\begin{tltxmdfexample}
  \surroundwithmdframed[linewidth=2pt]{verbatim}
\end{tltxmdfexample}

 
 \ExplCmd{mdflength\{选项\}} 
     若想使用 \Pack{mdframed} 定义的长度（例如
     \Opt{innerleftmargin}），可以使用命令 \Cmd{mdflength}。
\begin{tltxmdfexample}
  间距为 \hspace{\mdflength{innerleftmargin}} 
  \the\mdflength{innerleftmargin}
\end{tltxmdfexample}
 结果如下：
 \begin{center}
  间距为 \hspace{\mdflength{innerleftmargin}} \the\mdflength{innerleftmargin}
 \end{center}

 
 \ExplCmd{mdfsetup\{选项\}}设置选项时，
 既可以使用 \Cmd{usepackage} 的可选参数，
 也可以使用命令 \Cmd{mdfsetup}，后者不限于导言区。
 在一个分组内，设置仅作用于局部。\par
 \textbf{\color{red!70!black}在此我建议使用
 命令 \Cmd{mdfsetup}，而不是通过 \Cmd{usepackage} 的可选参数
 设置宏包选项，这样可以避免非健壮命令出错。}^^A
 \footnote{感谢 Heiko Oberdiek 与 Philipp Stephani
 \href{https://groups.google.com/group/de.comp.text.tex/^^A
       browse_thread/thread/15c5ad290490d426?hl=de}^^A
       {kvoptions 选项声明失败的讨论}}
 唯一的例外是选项 \Opt{framemethod}，它必须作为
 \Cmd{usepackage} 的可选参数给出。

 
 
 
 \section{定义自定义样式}
 下一节将介绍所有选项。不过，除了把选项作为全局选项或环境
 可选参数传入之外，你还可以定义自己的样式。方法如下：
 \ExplCmd{mdfdefinestyle\{样式名\}\{选项\}} 
   \Cmd{mdfdefinestyle} 允许用户定义不同的样式，再通过
   \Opt{style} 作为 \Pack{mdframed} 的选项使用。
 \ExplCmd{mdfapptodefinestyle\{样式名\}\{选项\}} 
      该命令用于向已定义的样式追加选项。\footnote{感谢
       Martin Scharrer 与 Enrico Gregorio：\newline
      \href{http://tex.stackexchange.com/questions/34684/argument-of-setkeys}
           {http://tex.stackexchange.com/questions/34684/argument-of-setkeys}}
 \ExplOpt{style} 
      若已用 \Cmd{mdfdefinestyle} 定义了某个样式，就可以用
      键 \Opt{style} 载入该样式。\Pack{mdframed} 目前
      没有预定义样式。

 下面是一个小例子：
\begin{tltxmdfexample}[morekeywords={mystyle}]
 \mdfdefinestyle{mystyle}{leftmargin=1cm,linecolor=blue}
 \begin{mdframed}[style=mystyle]
  文本
 \end{mdframed}
\end{tltxmdfexample}
 
 
 %%%%%%%%%%%%%%%%%%%%%%%
 \section{选项}
 本宏包提供多种用于调整边框的选项，下一节逐一列出。
 文档未列出部分可供调整的内部宏。
 所列选项分为全局选项与局部选项；全局选项
 不能在 \Cmd{mdfsetup} 内部使用。
 
 \subsection{全局选项}\label{globopt}
 以下选项仅为全局选项。
 
 \ExplOpt[none]{xcolor}
 设置该键后，宏包 \Pack{xcolor} 将以给定的
 参数被加载。若不提供任何值，\mbox{\Pack{mdframed}}
 会不带选项地加载宏包 \Pack{color}。若
 \Pack{xcolor} 已被加载，则给定的选项将被忽略。
 建议在 \Pack{mdframed} 之前加载 \Pack{xcolor}。
 
 \ExplOpt[default]{framemethod} 
     通过该键可以改变边框的绘制方式。可以选择边框由
 以下方式之一绘制：
 \begin{enumerate}
  \item \LaTeX 命令 {\small\Cmd{hrule}, \Cmd{vrule}, \Cmd{rule}},
  \item \Pack{TikZ} {\small（将加载宏包 \Pack{TikZ}）} 或
  \item \Pack{PSTricks} {\small（将加载宏包 \Pack{pstricks}）}。
 \end{enumerate}
 选项 \Opt{framemethod} 接受一个字符串，允许的组合
 列于下表。
 \begin{center}
 \renewcommand*\arraystretch{1.3}
 \captionof{table}{\Opt{framemethod} 的可用键}
 \begin{tabular}{@{}>{\small\bfseries}l>{\small\ttfamily}l@{}}
  \toprule
   \multicolumn{1}{@{}l}{\textbf{方式}} & \multicolumn{1}{l}{可用键}\\
  \midrule
   \LaTeX 命令   & default, tex, latex, none, 0  \\
   \Pack{TikZ}     & tikz, pgf, 1                  \\
   \Pack{PSTricks} & pstricks, ps, postscript, 2   \\
  \bottomrule
 \end{tabular}
 \end{center}

 图省事的话，也可以使用以下简写：
 \Opt{tikz}, \Opt{TikZ}, \Opt{pstricks}, \Opt{PSTricks} 或 \Opt{ps}。
 当然，默认方式没有对应的键。

 \minisec{备注}
 上述方式名无论写成大写、小写或大小写混合，
 效果都相同。

 下面列出的其余选项均可全局或局部设置，且不限于导言区。
 我尽量让选项名称做到见名知义。

 \subsection{恢复设置}
 有时需要清除此前所做的全部设置。每个默认值
 都保存在一个样式中，可通过选项 \Opt{default} 调用，
 或用更直观的名称
 \Opt{style=defaultoptions}；两者效果相同。

 下面是一个小例子：
\begin{tltxmdfexample}[morekeywords={mystyle}]
 \mdfsetup{%
   middlelinecolor=red,
   middlelinewidth=2pt,
   backgroundcolor=red!10,
   roundcorner=10pt}
 \begin{mdframed}
 文本\par 文本
 \end{mdframed}

 \begin{mdframed}[default]
 文本\par 文本
 \end{mdframed}
\end{tltxmdfexample}
 
 \begin{mdframed}[%
 backgroundcolor=yellow!70!brown,linewidth=1.5pt,
 roundcorner=8pt,linecolor=brown!70!black,innerleftmargin=10pt]
  \mdfsetup{default,
   middlelinecolor=red,
   middlelinewidth=2pt,
   backgroundcolor=red!10,
   roundcorner=10pt}
 \begin{mdframed}
 文本\par 文本
 \end{mdframed}

 \begin{mdframed}[default]
 文本\par 文本
 \end{mdframed}
 \end{mdframed}



 \subsection{长度选项}\label{optlength}
 
 图 \eqref{fig:laengen} 给出了可调整的长度
 （另见图 \eqref{fig:basicframes}）。
 
 \begin{center}
 \begin{tikzpicture}[scale=0.9]
      %Declaration
      \newlength{\linew}
      \setlength{\linew}{5pt}
      \tikzstyle{every node}=[font=\bfseries\large\sffamily]
      \coordinate (innersep) at (3.2,2.25);
      \coordinate (outersep) at (3.2,2);
      %Center
      \node[fill=green!10,draw=black,framed,rounded corners,
            minimum width=3cm,minimum height=1.5cm]%
            (CenterNode) {内容};
      
      %frames
      \draw[draw=red!70!black,line width=\linew] %
                ($(CenterNode)+(innersep)$) %
                rectangle ($(CenterNode)-(innersep)$);
      \draw[draw=blue!70!black,line width=\linew] %
                ($(CenterNode)+(innersep)+(\linew,\linew)$)%
                rectangle ($(CenterNode)-(innersep)+(-\linew,-\linew)$);
      \draw[draw=yellow!70!black,line width=\linew]
                ($(CenterNode)+(innersep)+(2\linew,2\linew)$)%
                rectangle ($(CenterNode)-(innersep)+(-2\linew,-2\linew)$);
      \draw[draw=black,line width=1] ($(CenterNode)+(innersep)+(outersep)$)%
                rectangle ($(CenterNode)-(innersep)-(outersep)$);
      
      
      \tikzstyle{every node}=[font=\small\ttfamily]
      %%innertopmargin
      \draw[draw=black,triangle 45-triangle 45]let \p1=(innersep) in
            (CenterNode.north)  -- ($(CenterNode)+(0,\y1)-(0,.5\linew)$) ;
      \path let \p1=(innersep) in
            (CenterNode.north)
            edge 
            node[align=center,framed,fill=white,draw=white]{innertopmargin} 
            ($(CenterNode)+(0,\y1)-(0,.5\linew)$) ;
      %%innerbottommargin
      \draw[draw=black,triangle 45-triangle 45]let \p1=(innersep) in
            (CenterNode.south)  -- ($(CenterNode)-(0,\y1)+(0,.5\linew)$) ;
      \path let \p1=(innersep) in
            (CenterNode.south)  
            edge 
            node[align=center,framed,fill=white,draw=white]{innerbottommargin}
            ($(CenterNode)-(0,\y1)+(0,.5\linew)$) ;
      
      %%innerleftmargin
      \draw[draw=black,triangle 45-triangle 45]let \p1=(innersep) in
            (CenterNode.west)  -- ($(CenterNode)-(\x1,0)+(.5\linew,0)$) ;
      \path let \p1=(innersep) in
            (CenterNode.west)  edge 
            node[align=center,rotate=90,framed,fill=white,draw=white]
                {innerleftmargin}  
            ($(CenterNode)-(\x1,0)+(.5\linew,0)$) ;
      %%innerrightmargin
      \draw[draw=black,triangle 45-triangle 45]let \p1=(innersep) in
            (CenterNode.east)  -- ($(CenterNode)+(\x1,0)-(.5\linew,0)$) ;
      \path let \p1=(innersep) in
            (CenterNode.east)  edge
            node[align=center,rotate=90,framed,fill=white,draw=white]
               {innerrightmargin}  
            ($(CenterNode)+(\x1,0)-(.5\linew,0)$) ;
      
      
      %%leftmargin
      \draw[draw=black,triangle 45-triangle 45]
            let \p1=(innersep) , \p2=(outersep) in
            ($(CenterNode)-(\x1,0)-(2.5\linew,0)$)  --
            ($(CenterNode)-(\x1,0)-(\x2,0)$) ;
      \path let \p1=(innersep) , \p2=(outersep) in
            ($(CenterNode)-(\x1,0)-(2.5\linew,0)$)  edge 
            node[align=center,rotate=0,framed,fill=white,draw=white]
                {leftmargin}  
            ($(CenterNode)-(\x1,0)-(\x2,0)$) ;
      %%rightmargin
      \draw[draw=black,triangle 45-triangle 45]
            let \p1=(innersep) , \p2=(outersep) in
            ($(CenterNode)+(\x1,0)+(2.5\linew,0)$)  --
            ($(CenterNode)+(\x1,0)+(\x2,0)$) ;
      \path let \p1=(innersep) , \p2=(outersep) in
            ($(CenterNode)+(\x1,0)+(2.5\linew,0)$)  edge 
            node[align=center,rotate=0,framed,fill=white,draw=white]
            {rightmargin}  
            ($(CenterNode)+(\x1,0)+(\x2,0)$)  ;
      
      %%skipabove
      \draw[draw=black,triangle 45-triangle 45]
            let \p1=(innersep) , \p2=(outersep) in
            ($(CenterNode)+(0,\y1)+(0,2.5\linew)$)  -- 
            ($(CenterNode)+(0,\y1)+(0,\y2)$) ;
      \path let \p1=(innersep) , \p2=(outersep) in
            ($(CenterNode)+(0,\y1)+(0,2.5\linew)$)  edge 
            node[align=center,rotate=0,framed,fill=white,draw=white]
                {skipabove}
            ($(CenterNode)+(0,\y1)+(0,\y2)$) ;
      %%skipbelow
      \draw[draw=black,triangle 45-triangle 45]
            let \p1=(innersep) , \p2=(outersep) in
            ($(CenterNode)-(0,\y1)-(0,2.5\linew)$)  -- 
            ($(CenterNode)-(0,\y1)-(0,\y2)$) ;
      \path let \p1=(innersep) , \p2=(outersep) in
            ($(CenterNode)-(0,\y1)-(0,2.5\linew)$)  edge 
            node[align=center,rotate=0,framed,fill=white,draw=white]
                {skipbelow}
            ($(CenterNode)-(0,\y1)-(0,\y2)$) ;
 \end{tikzpicture}
 
  \captionof{figure}{\Pack{mdframed} 可调长度}
  \label{fig:laengen}
 \end{center}
 
 \ExplOpt[pt]{defaultunit}
 所有长度接受两种输入：第一种是长度（如 2pt），
 第二种是数值（如 2），它将乘以
 \Opt{defaultunit}。为便于查阅，下表列出了所有长度选项及其
 默认值。选项名称已尽量做到见名知义，
 故不再逐一说明；但下面仍会介绍一些特殊的提示与长度。

 
 第一张表列出边框几何相关的长度。
 \begin{center}
 \def\arraystretch{1.4}
 \begin{tabular}{@{}%
 	  >{\Opt\bgroup}l<{\egroup}@{${}={}$}>{\ttfamily}l%
   @{\qquad}
 	  >{\Opt\bgroup}l<{\egroup}@{${}={}$}>{\ttfamily}l%
 	  @{}}
    \hline
 	  skipabove       & 0pt  & skipbelow         & 0pt  \\
 	  leftmargin      & 0pt  & rightmargin       & 0pt  \\
 	  innerleftmargin & 10pt & innerrightmargin  & 10pt \\
 	  innertopmargin  & 5pt  & innerbottommargin & 5pt  \\
    \hline
 \end{tabular}
 \end{center}
 在介绍下一组选项之前，先给出一个使用上述长度的例子。
\begin{tltxmdfexample}[morekeywords={mystyle}]
 \mdfdefinestyle{mdfexample1}{leftmargin=1cm,rightmargin=2cm,%
   innerleftmargin=1cm,innerrightmargin=1cm,roundcorner=10pt}
 \begin{mdframed}[style=mdfexample1]
   在任意直角三角形中，以斜边为边长的正方形面积，
   等于以两条直角边为边长的正方形面积之和。
 \end{mdframed}
\end{tltxmdfexample}
 \begin{mdframed}[%
 backgroundcolor=yellow!70!brown,linewidth=1.5pt,
 roundcorner=8pt,linecolor=brown!70!black,innerleftmargin=10pt]
 \mdfdefinestyle{mdfexample1}{default,leftmargin=1cm,rightmargin=2cm,%
   innerleftmargin=2cm,innerrightmargin=1cm,roundcorner=10pt}
 \begin{mdframed}[style=mdfexample1]
   在任意直角三角形中，以斜边为边长的正方形面积，
   等于以两条直角边为边长的正方形面积之和。
 \end{mdframed}
 \end{mdframed}


 下一张表列出用于调整 \Env{mdframed} 线条的选项。
 若使用 \Opt{framemethod=default}，边框只有一条线。
 \begin{center}
 \def\arraystretch{1.4}
 \begin{tabular}{@{}%
 	>{\Opt\bgroup}l<{\egroup}@{${}={}$}>{\ttfamily}l%
   @{\qquad}
 	>{\Opt\bgroup}l<{\egroup}@{${}={}$}>{\ttfamily}l%
   @{\qquad}
 	>{\Opt\bgroup}l<{\egroup}@{${}={}$}>{\ttfamily}l%
 	@{}}
     \hline
 	     linewidth       & 0.4pt  & innerlinewidth & 0pt  &
 	     middlelinewidth & 0.4pt \\
 	     outerlinewidth  & 0pt    & roundcorner     & 0pt &   
        \multicolumn{2}{c}{}    \\
      \hline
 \end{tabular}
 \end{center}
 若使用 \Opt{framemethod=tikz} 或 \Opt{framemethod=pstricks}，
 选项 \Opt{linewidth} 是选项 \Opt{middlelinewidth} 的别名。
 
 
 \subsection{颜色选项}
 下面给边框增添一些色彩。限制条件与长度选项相同：
 \Opt{framemethod=default} 下只有一条线。
 \begin{center}
 \def\arraystretch{1.4}
 \begin{tabular}{@{}%
 	>{\Opt\bgroup}l<{\egroup}@{${}={}$}>{\ttfamily}l%
   @{\qquad}
 	>{\Opt\bgroup}l<{\egroup}@{${}={}$}>{\ttfamily}l%
 	@{}}
     \hline
 	     linecolor       & black  & innerlinecolor  & black  \\
 	     middlelinecolor & black  & outerlinecolor  & black  \\
        backgroundcolor & white  & fontcolor       & black  \\
      \hline
 \end{tabular}
 \end{center}
 
 \ExplOpt[\{\}]{font} 
 虽说它并非真正的颜色选项，但你也可以指定内容的
 字体，这可通过选项 \Opt{font} 实现。

 掌握了长度和颜色之后，下面来改进之前的例子。
\begin{tltxmdfexample}[morekeywords={example1}]
 \mdfapptodefinestyle{example1}{backgroundcolor=brown!20,%
    linecolor=red!40!black,linewidth=4pt}
 \begin{mdframed}[style=mdfexample1]
   在任意直角三角形中，以斜边为边长的正方形面积，
   等于以两条直角边为边长的正方形面积之和。
 \end{mdframed}
\end{tltxmdfexample}
 \begin{mdframed}[%
 backgroundcolor=yellow!70!brown,linewidth=1.5pt,
 roundcorner=8pt,linecolor=brown!70!black,innerleftmargin=10pt]
 \mdfdefinestyle{mdfexample1}{default,leftmargin=1cm,rightmargin=2cm,%
   innerleftmargin=2cm,innerrightmargin=1cm,roundcorner=10pt,%
   backgroundcolor=brown!20,linecolor=red!40!black, linewidth=4pt}
 \begin{mdframed}[style=mdfexample1]
   在任意直角三角形中，以斜边为边长的正方形面积，
   等于以两条直角边为边长的正方形面积之和。
 \end{mdframed}
 \end{mdframed}



 \subsection{阴影}\label{sec:shadow}
 突出边框的另一种方式是添加阴影。\Pack{mdframed}
 提供了若干基础阴影选项。对于 \Opt{framemethod=tikz}，
 需要由用户通过 \verb+\usetikzlibrary{shadows}+ 加载
 \Pack{shadows} 库。
 \ExplOpt[8\,pt]{shadowsize}
 \ExplOpt[black!50]{shadowcolor}
 是否使用阴影取决于选项 \Opt{shadow}，取值可为
 \texttt{true} 或 \texttt{false}。启用阴影后会出现两个新选项：
 \Opt{shadowsize} 与 \Opt{shadowcolor}。
\begin{tltxmdfexample}[morekeywords={example1}]
 \begin{mdframed}[backgroundcolor=gray!40,shadow=true,roundcorner=8pt]
   在任意直角三角形中，以斜边为边长的正方形面积，
   等于以两条直角边为边长的正方形面积之和。
 \end{mdframed}
\end{tltxmdfexample}
 \begin{mdframed}[backgroundcolor=gray!40!white,shadow=true,roundcorner=8pt]
   在任意直角三角形中，以斜边为边长的正方形面积，
   等于以两条直角边为边长的正方形面积之和。
 \end{mdframed}


 \subsection{隐藏线条}
 前文图示展示了输出的默认行为。
 \Pack{mdframed} 允许只绘制指定的线条。若不想要右边线，
 可写 \Opt{rightline=false}。下表列出与此相关的所有键
 及其默认值。
 \begin{center}
 \def\arraystretch{1.4}
 \begin{tabular}{@{}%
 	>{\Opt\bgroup}l<{\egroup}@{${}={}$}>{\ttfamily}l%
   @{\qquad}
 	>{\Opt\bgroup}l<{\egroup}@{${}={}$}>{\ttfamily}l%
 	@{}}
     \hline
 	     topline  & true  & rightline  & true  \\
 	     leftline & true  & bottomline & true  \\
      \hline
 \end{tabular}
 \end{center}
 \ExplOpt[false]{hidealllines} 
 若想禁用全部线条，可以使用简写
 \Opt{hidealllines}。


 \subsection{双栏（twoside）模式}
 
 许多书籍采用 \texttt{twoside} 模式排版，此时
 最好按页面内外侧设置边距。除前面介绍的
 \Opt{leftmargin} 与 \Opt{rightmargin} 外，还可以使用
 选项 \Opt{outermargin} 与 \Opt{innermargin}。当然，在
 \texttt{oneside} 模式下这两个选项无效。即使文档采用
 \texttt{twoside} 模式，若不想使用 \Opt{outermargin} 与
 \Opt{innermargin}，也可通过选项
 \Opt{usetwoside} 关闭该行为。该键的取值为 \texttt{true} 或
 \texttt{false}。
 

 \subsection{脚注}
 在环境内部可以照常使用命令 \Cmd{footnote}。
 \Pack{mdframed} 采用 \Pack{minipage} 环境的语法，
 并共用同一套计数器。
 
 所有脚注文字都会收集到一个框中，显示在
 \Pack{mdframed} 环境的末尾。
 
 \ExplOpt[\mbox{} \Cmd{bigskipamount}]{footnotedistance} 
      该长度为 \Pack{mdframed} 环境末尾与
      \Cmd{footnoterule} 之间的距离。
 \ExplOpt[true]{footnoteinside} 
      脚注的位置可通过选项
      \Opt{footnoteinside} 改变。脚注显示在环境末尾，
      但可以决定输出位于 \Pack{mdframed} 框内
      还是框外。
 
 \minisec{注意}
 当选项 \Opt{footnoteinside=false} 时，脚注输出
 不处于被拆分的边框之内。我认为这样并不可取，因为
 新页面的第一行不宜是脚注。




\subsection{分页}
 \Pack{mdframed} 会在必要时拆分内容。图
 \ref{fig:basicframes} 给出了默认的拆分样式。不过
 在某些情况下不应发生分页。
 \ExplOpt[false]{nobreak} 
 为此可使用选项 \Opt{nobreak}，取值为
 \texttt{true} 或 \texttt{false}。如图
 \ref{fig:basicframes} 所示，不同部分的绘制方式不同，
 例如 \Pack{mdframed} 的中间部分只有左右两条边线。
 
 \ExplOpt[false]{everyline} 
 若希望 \Pack{mdframed} 的每一部分都拥有完整边框，可使用
 选项 \Opt{everyline}，此时各部分都按独立边框显示。

 前一节介绍了选项 \Opt{innertopmargin} 与
 \Opt{innerbottommaring}。发生分页时，还有两个新的长度选项
 会影响分页处的间距。
 \ExplOpt[0\,pt]{splittopskip} 
   设置环境被拆分后，上半部分底部预留的间距长度。
 \ExplOpt[0\,pt]{splitbottomskip} 
   设置环境被拆分后，下半部分顶部预留的间距长度。
 
 \subsection{边框标题}
 本节介绍与边框标题相关的所有选项，
 不再按其属性分类。
 
 \ExplOpt[none]{frametitle} 
      为环境设置标题。将标题作为环境选项传入，
      写法如 \Opt{frametitle=\{边框标题\}}。
 \ExplOpt[\Cmd{normalfont}\Cmd{bfseries}]{frametitlefont}
      设置 \Opt{frametitle} 的字体格式。
 \ExplOpt[\Cmd{raggedleft}]{frametitlealignment}
      设置 \Opt{frametitle} 的对齐方式。该选项必须通过 \Cmd{mdfsetup} 设置。
 \ExplOpt[false]{frametitlerule} 
      将此键设为 \Opt{true}，可在边框标题
      与正文之间加一条分隔线。
 \ExplOpt[.2\,pt]{frametitlerulewidth}
      设置 \Env{mdframed} 正文与标题之间
      分隔线的宽度。
 \ExplOpt[5\,pt]{frametitleaboveskip}
      设置边框标题与 \Env{mdframed} 上边距之间的间距。
 \ExplOpt[5\,pt]{frametitlebelowskip}
      设置边框标题与其分隔线之间的间距。
 \ExplOpt[backroundcolor]{frametitlebackgroundcolor}
      设置边框标题的背景颜色。
 \ExplOpt[false]{repeatframetitle}
      在每一段拆分出的边框上重复显示标题。


 下图展示使用选项 \Opt{frametitle} 时各长度的含义。
 
 \begin{center}
 \begin{tikzpicture}[scale=0.9]
      %Declaration
      \tikzstyle{every node}=[font=\bfseries\large\sffamily]
      \coordinate (innersep) at (3.2,2.25);
      \coordinate (outersep) at (3.2,2);
      %Center box
      \node[fill=green!10,draw=red!70!black,very thick,fill opacity=.2,
            text opacity=1,minimum width=9cm,minimum height=3cm,]%
            (CenterNode) {框内内容};
      %frametitle box
      \node[fill=yellow!70,draw=red!70!black,very thick,fill opacity=.2,
            text opacity=1,minimum width=9cm,minimum height=3cm,anchor=south]%
            at (CenterNode.north) (frametitle) {边框标题};
 % frametitlerule
     \draw[line width=3pt,blue!60!black] (frametitle.south west) 
                                      -- (frametitle.south east)
              node[pos=1,anchor=west,text width=4cm] 
                  {\Opt{frametitlerulewidth}\\\Opt{frametitlerule=true}};
 % frametitlebelowskip
     \draw[very thick,blue!60!black,latex-latex,shorten >=0.2cm] 
          (frametitle.south) -- (frametitle.center)
          node[pos=0.5,anchor=west] 
             {\Opt{ frametitlebelowskip}};
 % frametitleaboveskip
     \draw[very thick,blue!60!black,latex-latex,shorten >=0.2cm] 
          (frametitle.north) -- (frametitle.center)
           node[pos=0.5,anchor=west] {\Opt{ frametitleaboveskip}};
 % innertopmargin
     \draw[very thick,blue!60!black,latex-latex,shorten >=0.2cm] 
          (CenterNode.north) -- (CenterNode.center)
               node[pos=0.5,anchor=west] {\Opt{ innertopmargin}};
 % splittopskip
     \draw[very thick,blue!60!black,latex-latex,shorten >=0.2cm] 
          (CenterNode.north) -- (CenterNode.center)
               node[pos=0.5,anchor=east] {\Opt{splittopskip}};
 % global frame
  \draw[thick] (-5.5,-2.1) rectangle (9.4,5.5);
 \end{tikzpicture}
  \captionof{figure}{使用 \Opt{frametitle} 时各长度的含义}
  \label{fig:frametitlelaengen}
 \end{center}
 

 \minisec{备注与注意}
 边框标题的拆分是一个相当棘手的问题。若使用
 选项 \Opt{repeatframetitle}，拆分几乎必然出错；
 另一方面，使用该选项时用户必须妥善准备
 内容。

 \subsection{环境内部的标题命令}

 若要在 \Pack{mdframed} 环境内部提供小标题，可使用
 命令 \Cmd{mdfsubtitle}。相关选项列于下文。

 \ExplCmd{mdfsubtitle} 
      在 \Pack{mdframed} 内部设置内部一级小标题。

 该命令接受一个可选参数和一个
 必选参数。可选参数用于设置 \Pack{mdframed} 选项，
 且仅在局部生效；小标题的第二个参数也允许
 分段。

 \Cmd{mdfsubtitle} 内部不支持换行分段。

\begin{tltxmdfexample}
  \mdfsubtitle[<选项>]{小标题}
\end{tltxmdfexample}


 \subsubsection{与一级标题相关的选项}
 \ExplOpt[false]{subtitleaboveline}
     是否在小标题上方绘制一条线。
 \ExplOpt[false]{subtitlebelowline}
     是否在小标题下方绘制一条线。
 \ExplOpt[\string\normalfont\string\bfseries]{subtitlefont}
     设置小标题的字体。
 \ExplOpt[white]{subtitlebackgroundcolor}
     设置小标题在上下两条线之间的
     背景颜色。
 \ExplOpt[black]{subtitleabovelinecolor}
     设置上方线条的颜色。
 \ExplOpt[black]{subtitlebelowlinecolor}
     设置下方线条的颜色。
 \ExplOpt[0.8\,pt]{subtitleabovelinewidth}
     设置上方线条的宽度。
 \ExplOpt[0.6\,pt]{subtitlebelowlinewidth}
     设置下方线条的宽度。
 \ExplOpt[\string\baselineskip]{subtitleaboveskip}
     设置绘制上方线之前的间距。
 \ExplOpt[1.2\string\baselineskip]{subtitlebelowskip}
     设置绘制下方线之后的间距。
 \ExplOpt[0.5\string\baselineskip]{subtitleinneraboveskip}
     设置上方线与小标题本身之间的间距。
 \ExplOpt[0.5\string\baselineskip]{subtitleinnerbelowskip}
     设置小标题与下方线之间的间距。

 下面用一个示例演示其效果：

\newmdenv[%
 roundcorner=5pt,nobreak=true,
 subtitlebelowline=true,subtitleaboveline=true,
 subtitlebackgroundcolor=yellow!70!white,
 backgroundcolor=blue!20!white,
 frametitle={定理},frametitlerule=true,
 frametitlebackgroundcolor=yellow!70!white,
 ]{subtitleenv}

 \begin{subtitleenv}
   若干文字 \ldots
   \mdfsubtitle{注记}
   若干文字 \ldots
 \end{subtitleenv}

\begin{tltxmdfexample}
\newmdenv[%
 roundcorner=5pt,
 subtitlebelowline=true,subtitleaboveline=true,
 subtitlebackgroundcolor=yellow!70!white,
 backgroundcolor=blue!20!white,
 frametitle={定理},frametitlerule=true,
 frametitlebackgroundcolor=yellow!70!white,
 ]{subtitleenv}
\begin{subtitleenv}
  一些文本 \ldots
  \mdfsubtitle{注记}
  一些文本 \ldots
\end{subtitleenv}
\end{tltxmdfexample}


     
 \subsection{通用选项}\label{genopt}
 有些选项无法归入上述类别，集中列于本节。
 \ExplOpt[false]{ntheorem} 
      设置此布尔键之前，必须先加载
      宏包 \Pack{ntheorem}。启用该选项后，会把
      \Cmd{theorempreskipamount} 与
      \Cmd{theorempostskipamount} 设为 0\,pt，以避免额外的垂直间距。
 \ExplOpt[0\,pt]{needspace} 
      有时需要在边框被拆分之前预留一个最小高度，
      这时可使用 \Opt{needspace}。
      该选项接受一个长度，用于设定边框拆分前
      所需的最小高度。
 \ExplOpt[false]{ignorelastdescenders} 
      尝试忽略 \Env{mdframed} 最后一行的下伸字母。该
      思路来自 Tobias Weh，实现方案由
      Stefan Lemke 提供。参见
      \href{http://tex.stackexchange.com/questions/47584/how-to-make-^^A
            mdframed-ignore-descenders-in-last-line}^^A
           {如何让 mdframed 忽略最后一行的下伸字母}。
 此外，还有一些可供调整的长度。
 \ExplOpt[\textbackslash linewidth]{userdefinedwidth} 
      设置整个 \Pack{mdframed} 环境的宽度。
      该宽度包含线宽与内部边距，
      外部边距不计入其中。

 \ExplOpt[left]{align} 
      有时需要对齐环境本身。
      为此提供了选项 \Opt{align}，可取以下
      字符串：
      \begin{itemize}
       \item \Opt{left}，
       \item \Opt{right}，以及
       \item \Opt{center}。
      \end{itemize}
      \Opt{left} 或 \Opt{right} 对齐取决于给定的长度
      \Opt{leftmargin} 与 \Opt{rightmargin}。后文会给出一个示例
      来演示这段仍显笨拙的中文说明。

 请谨慎使用此选项。


 \subsection{TikZ 选项}\label{tikzopt}
 \mbox{}
 \ExplOpt[\{\}]{tikzsetting}
      通过该键可以向 \Cmd{tikzset} 传递若干选项。
      下一节给出了一些示例。务必
      将 \Opt{tikzsetting} 的选项置于花括号中。
 \ExplOpt[\{\}]{apptotikzsetting}
      通过该键可以向 \Opt{tikzsetting} 追加若干选项。
      该键基于对 \Pack{mdframed} 预定义键进行
      调整的思路。\Pack{mdframed} 通过
      \Cmd{tikzset} 定义了以下用于绘制边框的键。
      \begin{itemize}
       \item \Cmd{tikzset\{mdfbox/.style\}}
       \item \Cmd{tikzset\{mdfcorners/.style\}}
       \item \Cmd{tikzset\{mdfbackground/.style\}}
       \item \Cmd{tikzset\{mdfinnerline/.style\}}
       \item \Cmd{tikzset\{mdfouterline/.style\}}
       \item \Cmd{tikzset\{mdfmiddleline/.style\}}
       \item \Cmd{tikzset\{mdfframetitlerule/.style\}}
       \item \Cmd{tikzset\{mdfframetitlebackground/.style\}}
       \item \Cmd{tikzset\{mdfshadow/.style\}}
      \end{itemize}
      \textbf{修改其中任何一个之前，请先查看
              文件 \texttt{md-frame-1.mdf} 了解其设置。}


 \subsection{PSTricks 选项}\label{pstricksopt}
 \mbox{}
 \ExplOpt[\{\}]{pstrickssetting}
      通过该键可以向 \Cmd{psset} 传递若干选项。
      例如若希望所有线条为虚线，可设置
      \Opt{pstrickssetting=\{linestyle=dashed\}}。务必
      将 \Opt{pstrickssetting} 的选项置于花括号中。
 \ExplOpt[\{\}]{pstricksappsetting}
      \Pack{mdframed} 为不同元素使用预定义样式。结合本选项
      使用 \Cmd{addtopsstyle} 即可扩充这些定义。
      预定义样式包括：
      \begin{itemize}
       \item \Opt{mdfbackgroundstyle}
       \item \Opt{mdfframetitlebackgroundstyle}
       \item \Opt{mdfouterlinestyle}
       \item \Opt{mdfinnerlinestyle}
       \item \Opt{mdfmiddlelinestyle}
      \end{itemize}
       \textbf{修改其中任何一个之前，请先查看文件
               \texttt{md-frame-2.mdf} 了解其设置。}


 \section{钩子与布尔标志}\label{hooksandbools}
 以下选项与布尔标志可供各类用户使用；当然
 其中一些需要较深的使用经验。

 \ExplOpt[\{\}]{settings} 
 本选项允许用户在 \Pack{mdframed} 开头提交一些宏代码。
 所给内容在保存框开始之前执行。

 \ExplOpt[\{\}]{extra}
      向每一个绘制出的边框追加内容。
 \ExplOpt[\{\}]{singleextra}
      通过该键，可仅为未拆分的边框，向
      \Pack{mdframed} 的绘制环境追加额外内容。
 \ExplOpt[\{\}]{firstextra}
      通过该键，可仅为拆分边框的第一部分，向
      \Pack{mdframed} 的绘制环境追加额外内容。
 \ExplOpt[\{\}]{middleextra}
      通过该键，可仅为拆分边框的中间部分，向
      \Pack{mdframed} 的绘制环境追加额外内容，且仅作用于拆分边框的中间部分。 
      \limitation[TikZ、PSTricks]
 \ExplOpt[\{\}]{secondextra}
      通过该键，可仅为拆分边框的最后一部分，向
      \Pack{mdframed} 的绘制环境追加额外内容。


 在能向边框各部分追加内容之外，还可以判断
 当前绘制的是边框的哪一部分。为此定义了若干布尔标志，
 可用 \Cmd{ifbool\{布尔标志名\}} 检测。输出
 内容之前，所有标志均设为 true。
 \ExplOpt[false]{mdfsingleframe}
 该布尔仅在未拆分的边框上为 true。
 \ExplOpt[false]{mdffirstframe}
 该布尔仅在边框第一部分上为 true。
 \ExplOpt[false]{mdfmiddleframe}
 同上。
 \ExplOpt[false]{mdflastframe}
 同上。

 面向高级用户，还有一些可用于调整输出的钩子：
 \ExplOpt[\{\}]{beforesingleframe}
 传给本选项的代码在单个边框
 排版之前执行。
 \ExplOpt[\{\}]{aftersingleframe}
 传给本选项的代码在单个边框
 排版之后执行。
 \ExplOpt[\{\}]{beforebreak}
 本选项的值仅在可拆分的边框上执行。如前所述，
 可作用于第一段与中间段边框。
 \ExplOpt[\{\}]{afterbreak}
 本选项的值仅在可拆分的边框上执行。如前所述，
 可作用于第一段与中间段边框。
 \ExplOpt[\{\}]{beforelastframe}
 本选项仅在拆分边框的最后一段上执行。
 \ExplOpt[\{\}]{afterlastframe} 
 本选项仅在拆分边框的最后一段上执行。

 \ExplOpt[\{\}]{startcode}
 传给本选项的代码在环境
 开头执行。
 \ExplOpt[\{\}]{startinnercode}
 本选项在保存框开头执行。
 \ExplOpt[\{\}]{endinnercode}
 本选项在保存框末尾执行。
 \ExplOpt[\{\}]{endcode}
 本代码在环境末尾执行。

 \section{定理环境}
 本节介绍哪些命令可用于配合 \Pack{mdframed} 定义
 定理类环境。
 
 \ExplCmd{newmdtheoremenv} 
 由于本宏包常用来突出定理类环境，
 因而提供了一条简化该过程的命令。命令语法
 如下：
 
\begin{tltxmdfexample}[deletekeywords={mdframed}]
 \newmdtheoremenv[<mdframed 选项>]{<环境名>}%
                 [<编号参照>]{<标题>}[<编号层级>]
\end{tltxmdfexample}
 最后四个参数与命令 \Cmd{newtheorem} 相同。
 只有第一个可选参数用于传递 \Pack{mdframed} 选项。
 一个简单的示例如下：
\begin{tltxmdfexample}
 \theoremstyle{<定理样式>}
 \newmdtheoremenv[linecolor=blue]{lemma}{引理}[section]
 ...
 \begin{lemma}[标题示例]
   文本   文本   文本   文本   文本   文本
 \end{lemma}
\end{tltxmdfexample}
 目前还没有 \Cmd{renewmdtheoremenv}！


 \ExplCmd{mdtheorem} 
 这是 \Cmd{newtheorem} 的一种特殊形式。命令语法
 如下。
\begin{tltxmdfexample}[deletekeywords={mdframed}]
 \mdtheorem[<mdframed 选项>]{<环境名>}%
                 [<编号参照>]{<标题>}[<编号层级>]
\end{tltxmdfexample}
 可以看到，各参数与 \Cmd{newtheorem} 相同，但该命令
 会忽略所有 \Cmd{theoremstyle}。这是基于以下
 设计。

 命令 \Cmd{mdtheorem} 根据给定的第一个必选参数
 创建两个环境。第一个环境以该参数命名，生成带编号的定理；
 第二个环境在第一个必选参数名后加星号，
 排版格式相同，但不带编号。

 新定义环境的语法与普通定理
 环境相同。
\begin{tltxmdfexample}[deletekeywords={mdframed}]
 \begin{环境}[可选标题]
  ...
 \end{环境}
\end{tltxmdfexample}
 原理是：命令的标题文字会被用作边框标题。
 这样便可使用边框标题的全部选项。此外
 \Pack{mdframed} 还提供了下文说明的附加选项。

 \ExplOpt[\{:\}]{theoremseparator}
      设置定理编号与标题之间的分隔符。
      仅当给出定理标题时，才会打印
      \Opt{theoremseparator}。
 \ExplOpt[\{\}]{theoremtitlefont}
      通过选项 \Opt{frametitlefont} 可以调整
      边框标题的字体；而选项 \Opt{theoremtitlefont} 允许
      为定理标题设置另一种字体。
 \ExplOpt[\protect\Cmd{space}]{theoremspace}
      设置 \Opt{theoremseparator} 之后的间距。

 示例可参见随附文件。




 \section{复杂示例 —— Matlab 命令行窗口}
 下面这个示例改编自 TeX.SX 上的一个提问。
 \definecolor{DarkBlue}{rgb}{.11,.23,.60}
 \mdfdefinestyle{commandline}%
 {leftmargin=5pt, rightmargin=10pt,innerleftmargin=15pt,
  middlelinecolor=DarkBlue,
  middlelinewidth=2pt,
  frametitlerule=false,
  backgroundcolor=black!10!white,
  frametitle={命令窗口},
  frametitlefont={\normalfont\sffamily\color{white}\hspace{-1em}},
  frametitlebackgroundcolor=DarkBlue,
  singleextra={\draw[black!20,line width=12pt] 
       ($(O)+(7pt,1pt)$) --
       ($(O|-P)+(7pt,-\mdfframetitleboxtotalheight)-(0,1pt)$);
       \node[inner sep=0pt,color=black]at ($(O)+(7pt,9pt)$)%
       {$\scriptstyle f\!x$}; },
  nobreak=true,
 }
 
 \lstnewenvironment{script} {%
   \lstset{language=Matlab,basicstyle=\tiny\ttfamily,breaklines=true,%
           aboveskip=0pt,belowskip=0pt}}{}
 \surroundwithmdframed[style=commandline]{script}
 \begin{script}
 >> help sin
  sin    以弧度为参数的正弦函数。
     sin(X) 表示 X 中各元素的正弦值。
 
     另见 asin、sind。
 
     重载方法：
        sdpvar/sin
        codistributed/sin
        gpuArray/sin
 
     帮助浏览器中的参考页
        doc sin
 
 >> 
 \end{script}

 得到该结果的代码如下：
\begin{tltxmdfexample}[moretexcs={lstnewenvironment,definecolor,lstset}]
 \definecolor{DarkBlue}{rgb}{.11,.23,.60}
 \mdfdefinestyle{commandline}%
 {leftmargin=5pt, rightmargin=10pt,innerleftmargin=15pt,
  middlelinecolor=DarkBlue,
  middlelinewidth=2pt,
  frametitlerule=false,
  backgroundcolor=black!10!white,
  frametitle={命令窗口},
  frametitlefont={\normalfont\sffamily\color{white}\hspace{-1em}},
  frametitlebackgroundcolor=DarkBlue,
  singleextra={\draw[black!20,line width=12pt] 
       ($(O)+(7pt,1pt)$) --
       ($(O|-P)+(7pt,-\mdfframetitleboxtotalheight)-(0,1pt)$);
       \node[inner sep=0pt,color=black]at ($(O)+(7pt,9pt)$)%
       {$\scriptstyle f\!x$}; },
  nobreak,
 }
 
 \lstnewenvironment{script} {%
   \lstset{language=Matlab,basicstyle=\tiny\ttfamily,breaklines=true,%
           aboveskip=0pt,belowskip=0pt}}{}
 \surroundwithmdframed[style=commandline]{script}
 \begin{script}
 >> help sin
  sin    以弧度为参数的正弦函数。
     sin(X) 表示 X 中各元素的正弦值。
 
     另见 asin、sind。
 
     重载方法：
        sdpvar/sin
        codistributed/sin
        gpuArray/sin
 
     帮助浏览器中的参考页
        doc sin
 
 >> 
 \end{script}
\end{tltxmdfexample}

 
 
 
 \section{示例}
 
 我把示例拆分到四个文档文件中。这些文件
 分别是
 \begin{description}
 \item[\texttt{mdframed-example-default}] 
       演示用 \Opt{framemethod=default} 创建的示例。
 \item[\texttt{mdframed-example-tikz}]
       演示用 \Opt{framemethod=TikZ} 创建的示例。
 \item[\texttt{mdframed-example-pstricks}] 
       演示用 \Opt{framemethod=pstricks} 创建的示例。
 \item[\texttt{mdframed-example-texsx}] 
       演示与 \Pack{listings} 等交互的示例。
 \end{description}
 
 各示例通常并不完全等同，但一般可以改用到
 其他绘制方式上。
 因此我强烈建议翻阅所有示例文件。

 韩国 \TeX 小组制作了一份非常精彩的演示文稿。我想分享这个
 链接，因为它确实是一件很棒的作品：
 \href{http://prezi.com/zc355i9tcsgz/talk-kts-2012-mdframed/}^^A
      {kts 2012 mdframed}。
 

 \IfFileExists{mdframed-example-default.pdf}%
   {\includepdf[pages=-]{mdframed-example-default.pdf}}{}
 \IfFileExists{mdframed-example-tikz.pdf}%
   {\includepdf[pages=-]{mdframed-example-tikz.pdf}}{}
 \IfFileExists{mdframed-example-texsx.pdf}%
   {\includepdf[pages=-]{mdframed-example-texsx.pdf}}{}
 \IfFileExists{mdframed-example-pstricks.pdf}%
   {\includepdf[pages=-]{mdframed-example-pstricks.pdf}}{}

 \section{错误、警告与提示信息}
 
 \Pack{mdframed} 会在 \texttt{日志} 文件中输出不同的错误、警告与
 提示信息。有些 \LaTeX 编辑器，如
 \TeX Maker 或 \TeX Studio，设有专门的错误与警告标签页，
 但不显示提示信息，因此应直接查看 \texttt{日志文件}。
 
 以下为 \Pack{mdframed} 产生的错误与警告。
 
\begin{tltxmdfexample}
   宏包 ... 不存在，但 mdframed 需要它
\end{tltxmdfexample}
 要避免此问题，应安装
 \ref{sec:syntax} 节列出的所需宏包。
 
\begin{tltxmdfexample}
   宏包选项 style 已弃用
   请改用 framemethod 代替 style
\end{tltxmdfexample}
 自 0.9d 版本起，\Pack{mdframed} 改变了选项
 \Opt{style} 的含义：该选项用于载入由
 \Cmd{mdfdefinestyle} 定义的样式。请改用 \Opt{framemethod}
 （见 \ref{globopt} 节）。
 
\begin{tltxmdfexample}
  未知的 framemethod .... mdframed
\end{tltxmdfexample}
 选项 \Opt{framemethod} 的输入字符串无法识别。参见
 \ref{globopt} 节。
 
\begin{tltxmdfexample}
   尚未加载 ntheorem
\end{tltxmdfexample}
 要使用选项 \Opt{ntheorem}，必须先
 加载宏包 \Pack{ntheorem}。
 
\begin{tltxmdfexample}
  内容宽度仅有 3cm
\end{tltxmdfexample}
 \Pack{mdframed} 根据给定选项计算内容宽度。
 若内容宽度小于
 3\,cm，就会出现此警告。应调整
 设置以获得更大的宽度。
 
\begin{tltxmdfexample}
  出现了不良分页
  必须手动调整：修改文字、间距
  或其他内容
\end{tltxmdfexample}
 有时垂直空间足以容纳边线以及边线与内容之间的间距，
 却不足以容纳内容本身。
 此时会给出此警告，因为该框的
 内容为空。可以调整设置，或在
 \Pack{mdframed} 环境之前插入 \Cmd{clearpage}。
 目前我还没有办法彻底避免这类情况。
 
\begin{tltxmdfexample}
  出现了不良分页
  因为被拆分的框为空
  必须调整页面设置，例如
  使用 enlargethispage 等手段
  出现了不良分页
\end{tltxmdfexample}
 理由同上。
 
\begin{tltxmdfexample}
  出现了不良分页
  因为最后一个被拆分的框为空
  必须调整设置
\end{tltxmdfexample}
 理由同上，但仅出现在最后一个框中。
 
 
\begin{tltxmdfexample}
  选项 ... 已被占用
  在第 ... 行输入中不再生效
\end{tltxmdfexample}
 若在文档
 正文中设置全局选项，就会出现此警告。
 
 \section{已知局限}
  本节汇总已知问题。若遇到
  其他问题，请发邮件至
  \href{mailto:marco.daniel@mada-nada.de}{marco.daniel at mada-nada.de}。

    对本宏包的进一步扩展，你有什么想法或期望吗？
    欢迎告诉我！
 
 \begin{enumerate}
  \item 目前该环境与宏包
        \Pack{gmverb} 不兼容。
  \item 若加载宏包 \Pack{picins}，边框将无法拆分。
        这源于宏包 picins 的一个问题：它把
        \Cmd{@captype} 定义为全局。
        要配合 picins 使用，可采用下面的临时办法。
\begin{tltxmdfexample}
\usepackage{picins}
\makeatletter
\let\@captype\@undefined
\def\newcaption{%
 \begingroup%
  \def\@captype{figure}%
  \refstepcounter\@captype\@dblarg{\@newcaption\@captype}%
 \endgroup%
}
\makeatother
\end{tltxmdfexample}
  \item \Pack{mdframed} 无法处理
        文档类 \Pack{beamer} 的选项 \Opt{allowframebreaks}。
 \item 嵌套的 \Env{mdframed} 环境无法拆分。 
 \end{enumerate}
 
 
 \section{待办事项}
 \textbf{\color{red} 更新文档十分重要}
 \begin{enumerate}
   \item 见“已知问题”。
   \item 目前还无法把
         宏包 \Pack{multicol} 的环境
         \Cmd{begin\{multicols\}} 与 \Pack{mdframed} 的全部选项配合使用。
   \item 创建新的样式。
   \item 改进分页。
   \item 改进脚注。
   \item 改进文档与示例。
   \item 为 \Opt{frametitle} 创建样式。
   \item 提供 \Env{mdframed} 的行内版本，行为类似
         \Cmd{fbox}。
   \item 在文件 \Pack{md-frame-1.mdf} 中加入 \Cmd{ht}\Cmd{strutbox}。
 \end{enumerate}                 
 
 \section{致谢}
  \textsf{Dick Nickalls;\quad Dietrich Grau;\quad Piazza Luca;\quad%
          Jobst Hoffmann;\quad Martin Scharrer;\quad Enrico Gregorio;\quad 
          Heiko Oberdiek;\quad Philipp Stephani}。
 \medskip
 
 \noindent 感谢校对
 
 \textsf{Alan Munn} 与 \textsf{Nahid Shajari}
 
 希望没有遗漏任何人。
 \clearpage
 \appendix

 \section{更多信息}
 %\addtocontents{toc}{\protect\setcounter{tocdepth}{-4}}
 下一节介绍如何创建
 你自己的边框。
 \subsection{\Pack{mdframed} 的工作原理}
 通过环境 \Cmd{begin\{mdframed\}} \ldots \Cmd{end\{mdframed\}}，
 全部内容会保存在名为
 \Cmd{mdf@splitbox@one} 的 \Cmd{savebox} 中。在计算出
 \Cmd{mdf@splitbox@one} 的宽度与高度（由 \Pack{mdframed.sty} 完成）之后，
 依次放置该框（由 \Pack{md-frame-X.mdf} 完成）。下图演示了
 这一过程。
 
 \begin{center}
 \begin{tikzpicture}
 \newcommand{\mybox}[1]{\parbox{3.5cm}{\centering\small #1}}
 \tikzstyle{every node}=[font=\sffamily]
 \tikzset{mymatrixnodes/.style={%
      matrix of nodes,
      nodes={framed,draw=black,fill=yellow!10,},
      column sep=4.5cm,
      row sep=1.5cm,
      }}
 \matrix(m)[mymatrixnodes]
   {%
    \mybox{读入内容并存入 \Cmd{mdf@splitbox@one}}              &
    \mybox{\Cmd{mdf@putbox@single}}                                         \\
    \mybox{将 \textbackslash mdf@splitbox@one 拆分到
           页面剩余空间中的 \Cmd{mdf@splitbox@two}}                                       &
    \mybox{\Cmd{mdf@putbox@first}}                                          \\
    \mybox{\Cmd{mdf@splitbox@one} 的内容在新
           页面上有足够空间}                                                             &
    \mybox{\Cmd{mdf@putbox@second}}                                         \\
    \mybox{将 \Cmd{mdf@splitbox@one} 拆分到页面剩余空间中的
           \Cmd{mdf@splitbox@two}}                                                &
    \mybox{\Cmd{mdf@putbox@middle}}                                         \\
 };
 
 \tikzstyle{every path}=[black,-triangle 45,thick,shorten >=4pt,shorten <=4pt]
 \draw  (m-1-1) edge 
        node[align=center,anchor=south]{页面上空间足够} (m-1-2);
 \draw  (m-2-1) edge 
        node[align=center,anchor=south]{排出 \Cmd{mdf@splitbox@two}}
         (m-2-2);
 \draw  (m-3-1) edge 
        node[align=center,anchor=south]{是} (m-3-2);
 \draw  (m-4-1) edge 
        node[align=center,anchor=south]{排出 \Cmd{mdf@splitbox@two}}
         (m-4-2);
 
 \draw  (m-2-2) edge node[align=center,fill=white]%
        {计算新 \Cmd{mdf@splitbox@one} 的高度} (m-3-1);
 
 
 \draw  (m-1-1) edge 
        node[align=center,anchor=west]{页面上空间不足} (m-2-1);
 \draw  (m-3-1) edge node[align=center,anchor=west]{否} (m-4-1);
 
 \draw  let \p1=($(m-4-2.south)-(0,1)$), \p2=($(m-3-1.west)-(0.8,0)$) in
       (m-4-2.south) -- ($(m-4-2.south)+(0,-1)$) --
            node[align=center,anchor=north]%
            {计算新 \Cmd{mdf@splitbox@one} 的高度} (\x2,\y1)
        -- ($(m-3-1.west)-(0.8,0)$)-- ($(m-3-1.west)$) ;
 
 
 \end{tikzpicture}
 
 \captionof{figure}{\Pack{mdframed} 内容的安放方式}\label{fig:way}
 \end{center}
 
 内容宽度由以下设置共同决定：
 \Opt{leftmargin}、\Opt{rightmargin}、\Opt{linewidth}、
 \Opt{innerleftmargin} 与 \Opt{innerrightmargin}
 （见图 \eqref{fig:laengen}）。
 
 \subsection{边框命令}
 \Pack{mdframed} 认识四种“边框命令”。
 这些命令告诉 \LaTeX\ 如何安放 \Pack{mdframed} 的内容。 
 \begin{description}
   \item[\Cmd{mdf@putbox@single}] 
        安放单个未拆分边框的内容。
   \item[\Cmd{mdf@putbox@first}]  
        安放拆分边框第一段的内容。
   \item[\Cmd{mdf@putbox@middle}] 
        安放拆分边框中间段的内容。
   \item[\Cmd{mdf@putbox@second}] 
        安放拆分边框最后一段的内容。
 \end{description}
 
 借助上述命令给出一个示例。命令 \Cmd{box}
 取出保存框中的内容并排版输出。
 
 首先排出不带任何设置（但带有已计算宽度）的单个框。
\begin{tltxmdfexample}
 \makeatletter
  \def\mdf@putbox@single{\box\mdf@splitbox@one}
 \makeatother
\end{tltxmdfexample}
 这里使用命令 \Cmd{leftline}，使“边框命令”从
 左侧开始。
\begin{tltxmdfexample}[deletekeywords={leftline},moretexcs={leftline}]
 \makeatletter
  \def\mdf@putbox@single{\leftline{\box\mdf@splitbox@one}}
 \makeatother
\end{tltxmdfexample}
 
 接下来需要了解各长度的命名方式。所有可由选项
 修改的长度均采用以下语法：
\begin{tltxmdfexample}
 \mdf@<长度名称>@length
\end{tltxmdfexample}
 例如 leftmargin 对应的长度为：
\begin{tltxmdfexample}
 \mdf@leftmargin@length
\end{tltxmdfexample}
 
 若只想绘制一条带正确 \Opt{leftmargin} 的左边线，
 可按如下方式定义 \Cmd{mdf@putboxsingle}。
 
\begin{tltxmdfexample}[deletekeywords={leftline},moretexcs={leftline}]
 \makeatletter
 \def\mdf@putbox@single{%
          \leftline{%
              \hspace*{\mdf@leftmargin@length}%
              \rule[-\dp\mdf@splitbox@one]{\mdf@linewidth}%
                   {\ht\mdf@splitbox@one+\dp\mdf@splitbox@one}%
              \box\mdf@splitbox@one
          }%
 }
 \makeatother
\end{tltxmdfexample}
 
 这样便可实现你想要的效果。若创建自己的样式，
 可将文件保存为 \Pack{md-frame-X.mdf}。
 其中 \Pack{X} 必须为整数。
 之后即可通过选项 \Opt{framemethod}，
 设为 \Opt{framemethod=X} 来载入该文件。
 

\PrintIndex

\end{document}
